\section{Even quadratic recurrences and height asymptotics}
\label{sec:recurrences}

The modular criterion admits a complete elementary parametrization.
It also gives a way to enumerate the exceptional rational set without
factoring either entry.

\begin{theorem}[Recurrence parametrization]\label{thm:recurrences}
Apart from $J(1)$, all unordered reduced logarithmic pairs $1\le a<b$
are exactly the adjacent pairs in the following families, for every even
integer $c\ge2$:
\begin{align}
 u_0=1,\quad u_1=c,\quad u_{j+1}&=c u_j-u_{j-1},\label{eq:minus-rec}\\
 v_0=1,\quad v_1=c,\quad v_{j+1}&=c v_j+v_{j-1}.\label{eq:plus-rec}
\end{align}
Both $a/b$ and $b/a$ are logarithmic cases. The seed $(1,c)$ is common
to the two families. Beyond the seeds there is no overlap between the
families or between different values of $c$.
\end{theorem}
\begin{proof}
Let $e$ be even and $o$ odd in a pair satisfying
\eqref{eq:mixed-class}. First suppose $e^2\equiv1\pmod o$.
Then
\[
 c=\frac{e^2+o^2-1}{eo}
\]
is an integer. It is even: after division of the numerator by $e$, both
$e$ and $(o^2-1)/e$ are even, and the remaining divisor $o$ is odd.
Also $c\ge2$, since $(e-o)^2\ge1$. Writing the sorted pair as $a<b$
gives
\begin{equation}\label{eq:minus-invariant}
 a^2+b^2-cab=1.
\end{equation}
If $a>1$, the new entry
$d=(a^2-1)/b=ca-b$ is a positive integer strictly smaller than $a$.
The pair $(d,a)$ has the same invariant, is coprime, and has opposite
parity because $c$ is even and $d$ has the parity of $b$. The invariant
itself gives both congruences in \eqref{eq:mixed-class}. Repetition
reaches a pair $(1,c)$. Reversing the descent is
\eqref{eq:minus-rec}.

Next suppose $e^2\equiv-1\pmod o$. The signed integer
\[
 H=\frac{o^2-e^2-1}{eo}
\]
is even by the same divisibility argument, and is nonzero: otherwise
$(o-e)(o+e)=1$, impossible for $e>0$. Put $c=|H|\ge2$. In sorted
coordinates,
\begin{equation}\label{eq:plus-invariant}
 b^2-cab-a^2=\varepsilon,
 \qquad \varepsilon=\begin{cases}1,&b\text{ odd},\\-1,&b\text{ even}.
 \end{cases}
\end{equation}
For $a>1$ set
$d=b-ca=(a^2+\varepsilon)/b$. This is a positive integer less than
$a$; for $\varepsilon=1$, use $a^2+1<a(a+1)\le ab$. The pair
$(d,a)$ is coprime and has opposite parity. Its invariant has sign
$-\varepsilon$, exactly as required when the parity of the larger entry
changes. Thus the descent continues to $(1,c)$, and its reverse is
\eqref{eq:plus-rec}.

Conversely the recurrences preserve \eqref{eq:minus-invariant} or
\eqref{eq:plus-invariant}, respectively, and preserve coprimality and
alternating parity. In the minus family the invariant gives
$e^2\equiv1\pmod o$ and
$o^2-1=e(co-e)$, which is divisible by $2e$. In the plus family it
gives $e^2\equiv-1\pmod o$ and
$o^2-1=e(e\pm co)$, again divisible by $2e$. Thus every generated
pair satisfies the criterion.

If the odd entry exceeds $1$, the signs $+1$ and $-1$ modulo that
entry cannot both hold. The displayed invariants determine $c$
uniquely within each case, and the strictly decreasing descent
identifies a unique depth. The only possible shared pairs are the seeds,
whose odd entry is $1$.
\end{proof}

\begin{center}
\begin{tabular}{lll}
\toprule
Family & Coefficient & Initial terms\\
\midrule
Minus & $2$ & $1,2,3,4,5,6,\ldots$\\
Minus & $4$ & $1,4,15,56,209,\ldots$\\
Minus & $6$ & $1,6,35,204,1189,\ldots$\\
Plus & $2$ & $1,2,5,12,29,70,\ldots$\\
Plus & $4$ & $1,4,17,72,305,\ldots$\\
Plus & $6$ & $1,6,37,228,1405,\ldots$\\
\bottomrule
\end{tabular}
\end{center}

The coefficient-$2$ minus family explains why every neighboring ratio
is logarithmic. All other families grow geometrically. In particular,
a descent algorithm can shortcut consecutive integers and otherwise
uses $O(\log\max(a,b))$ steps. This statement counts arithmetic descent
steps; it is not a bound on the bit complexity of an arbitrary expanded
analytic certificate.

\subsection{A complete hierarchy of height terms}

Let $N(B)$ count unordered coprime pairs $1\le a<b\le B$ satisfying the
logarithmic criterion; $J(1)$ is not counted. Here $B$ is a positive
integer.

\begin{theorem}[Height expansion to every fixed order]\label{thm:counting}
For every fixed integer $R\ge2$, as $B\to\infty$,
\begin{equation}\label{eq:count-expansion}
 N(B)=\frac32 B+\sum_{r=2}^{R}B^{1/r}
       +O_R\bigl(B^{1/(R+1)}\log B\bigr).
\end{equation}
In particular,
\begin{equation}\label{eq:count-two}
 N(B)=\frac32 B+\sqrt B+O(B^{1/3}\log B).
\end{equation}
The number of ordered positive rational logarithmic arguments of height
at most $B$ is $2N(B)+1$. For $B\ge5$, the number admitting a formal
rational-dilogarithmic reduction is $2N(B)+7$.
\end{theorem}
\begin{proof}
Count each seed $(1,c)$ once: there are $\lfloor B/2\rfloor$.
The minus family at coefficient $2$ contributes all consecutive pairs;
after removing its seed it contributes $B-2$ for $B\ge2$. Their sum
is $\tfrac32B+O(1)$.

For the other pairs write $w_r^{\pm}(c)$ for the term at index $r$ in
the two recurrences. The pair with this height has indices $(r-1,r)$.
At fixed $r\ge2$, induction on the recurrence gives
\[
 w_r^{\pm}(c)=c^r\pm(r-1)c^{r-2}+O_r(c^{r-4}),
\]
where the remainder can simply be read as the lower-degree polynomial
terms; when $r$ is small it is zero as appropriate. Equivalently,
$w_r^{\pm}(c)=c^r+O_r(c^{r-2})$. The number of even coefficients
satisfying $w_r^{\pm}(c)\le B$ is therefore
$\tfrac12 B^{1/r}+O_r(1)$: the transition can differ from
$c=B^{1/r}$ by only $O_r(B^{-1/r})$, with finitely many small
coefficients absorbed in the constant. Removing coefficient $2$ from
the minus family changes this by $O(1)$. Both signs together contribute
$B^{1/r}+O_r(1)$. Summing over $2\le r\le R$ gives the displayed
fractional powers.

For completeness, the remaining depths are uniformly bounded. In the
minus family with $c\ge4$, induction on the successive ratios gives
$w_j/w_{j-1}>c-1$, so $w_j\ge(c-1)^j$. In the plus family with
$c\ge2$, $w_j\ge c^j$. A term with index at least $R+1$ and height
at most $B$ therefore requires
$c\le 1+B^{1/(R+1)}$. There are
$O(B^{1/(R+1)})$ possible coefficients, and each contributes only
$O(\log B)$ depths, because its growth factor is at least $2$.
This proves the remainder estimate. Disjointness from
Theorem~\ref{thm:recurrences} prevents overcounting.

Finally each unordered pair yields two reciprocal arguments, while
$x=1$ adds one. The six distinct extra arguments in
\eqref{eq:six-extras} all have height at most $5$.
\end{proof}

The theorem implies a sparse set among all reduced rational arguments
of bounded height, but no transcendence statement is attached to that
sparsity. It counts exactly the formal reduction cases.

\begin{center}
\begin{tabular}{r r r}
\toprule
$B$ & $N(B)$ & Ordered logarithmic arguments\\
\midrule
$100$ & $161$ & $323$\\
$1\,000$ & $1\,543$ & $3\,087$\\
$10\,000$ & $15\,135$ & $30\,271$\\
$1\,000\,000$ & $1\,501\,162$ & $3\,002\,325$\\
$100\,000\,000$ & $150\,010\,658$ & $300\,021\,317$\\
$10\,000\,000\,000$ & $15\,000\,102\,700$ & $30\,000\,205\,401$\\
\bottomrule
\end{tabular}
\end{center}
These are exact counts from the recurrence enumerator, not floating-point
fits to \eqref{eq:count-expansion}. The script uses integer arithmetic
and does not enumerate all $O(B^2)$ pairs at the larger heights.
